\documentclass{ctexbeamer}
\usepackage[UTF8]{ctex} % 添加此行确保 listings 支持 UTF8
\usepackage{graphicx}
\usepackage{multicol}

\usefonttheme{serif}
\usefonttheme{professionalfonts}
\definecolor{tyblue}{RGB}{102,204,255}
\usepackage{amsmath}
\usepackage{mathtools}
\usepackage{hyperref}
\usepackage{tikz}
\usetikzlibrary{positioning, shapes.geometric}

\usepackage{listings}
\usepackage{xcolor}

% 自定义代码块样式
\lstset{
    numbers=left,
    numberstyle=\tiny,
    basicstyle=\ttfamily\small, % 推荐加 \small
    keywordstyle=\color{blue!70},
    commentstyle=\color{red!50!green!50!blue!50},
    frame=shadowbox,
    rulesepcolor=\color{red!20!green!20!blue!20},
    escapeinside=``,
    xleftmargin=2em, xrightmargin=2em, aboveskip=1em,
    framexleftmargin=2em,
    language=C,
    extendedchars=true, % 支持扩展字符
    inputencoding=utf8 % 明确指定编码
} 
\usetheme{Antibes}
\usecolortheme{dolphin}

% 自定义 x86-64 汇编语言以改进关键字高亮
\lstdefinelanguage{x86_64}{
  morekeywords={
    mov,add,sub,and,or,xor,not,neg,cmp,test,
    je,jne,jl,jle,jg,jge,jmp,call,ret,
    shr,shl,sar,sal,lea,inc,dec,push,pop,pushq,movq,andl,
    sete,setne,cmove,cmovne,cdqe,cqo,movl,movslq,testl,js,subl,addl,cmpl
  },
  sensitive=true,
  morecomment=[l]{;},
  morestring=[b]"
}

\begin{document}

\title{回课: 控制流}
\subtitle{机器级编程 II}
\author{Rainybunny}
\date{\today}
\frame{\titlepage}

\begin{frame}{目录}
%   \begin{multicols}{2}
    \tableofcontents[hideallsubsections]
%   \end{multicols}
\end{frame}

\AtBeginSection[]{
  \begin{frame} % [shrink=15]
    % \begin{multicols}{2}
      \tableofcontents[currentsection,hideallsubsections]
    % \end{multicols}
  \end{frame}
}

\section{C loops 与 goto loops 的互译}

\subsection*{do-while loop}

\begin{frame}[fragile]{do-while loop}
  \begin{lstlisting}[language=C]
  do {
    count += x & 1;
    x >>= 1;
  } while (x);
  \end{lstlisting}

  $$
    \Updownarrow
  $$

  \begin{lstlisting}[language=C]
loop:
  count += x & 1;
  x >>= 1;
  if (x) goto loop;
  \end{lstlisting}
\end{frame}

\subsection*{while loop}

\begin{frame}[fragile]{while loop}
  \begin{lstlisting}[language=C]
  while (x) {
    count += x & 1;
    x >>= 1;
  }
  \end{lstlisting}

  $$
    \Updownarrow
  $$

  \qquad 方案一: jump to middle.
  \begin{lstlisting}[language=C]
  goto check;
loop:
  count += x & 1;
  x >>= 1;
check:
  if (x) goto loop;
  \end{lstlisting}
\end{frame}

\begin{frame}[fragile]{while loop}
  \begin{lstlisting}[language=C]
  while (x) {
    count += x & 1;
    x >>= 1;
  }
  \end{lstlisting}

  $$
    \Updownarrow
  $$

  \qquad 方案二: guarded-do.
  \begin{lstlisting}[language=C]
  if (!x) goto finished;
loop:
  count += x & 1;
  x >>= 1;
  if (x) goto loop;
finished: ;
  \end{lstlisting}
\end{frame}

\begin{frame}{while loop}
  \qquad 在不同优化级别下, \texttt{gcc} 对 \texttt{while} 循环的处理方式不完全相同. 在实验\footnote{https://godbolt.org/z/7KKMWs3Mz}中:
  \begin{itemize}
    \item 启用 \texttt{-Og} 选项, \texttt{while} 被处理为 jump-to-middle 方案;
    \item 启用 \texttt{-O1} 选项, \texttt{while} 被处理为 guarded-do 方案.
  \end{itemize}
  相比 guarded-do, jump-to-middle 会多一次跳转 (第一次成功进入循环体之前会连续跳转两次), 有轻微性能劣势.
\end{frame}

\subsection*{for loop}

\begin{frame}[fragile]{for loop}
  \begin{lstlisting}[language=C]
  for (int `{\color{red}\texttt{\small{i = 0}}}`; `{\color{green}\texttt{\small{x >> i}}}`; `{\color{cyan}\texttt{\small{++i}}}`) {
    count += x >> i & 1;
  }
  \end{lstlisting}

  $$
    \Updownarrow
  $$

  \begin{lstlisting}[language=C]
  int `{\color{red}\texttt{\small{i = 0}}}`;
  if (!(x `\underline{\texttt{\small{>> i}}}`)) goto finished;
loop:
  count += x >> i & 1;
  `{\color{cyan}\texttt{\small{++i;}}}`
  if (`{\color{green}\texttt{\small{x >> i}}}`) goto loop;
finished: ;
  \end{lstlisting}
\end{frame}

\subsection*{多重循环}

\begin{frame}[fragile, shrink=35]{多重循环}
  \begin{lstlisting}[language=x86_64]
  movl  $1000000000, 96(%rsp)
  movl  $9, %esi
  movl  $0, %eax
  jmp .L4
.L6:
  subl  $1, %eax
.L5:
  movslq  %eax, %rdx
  cmpl  %ecx, 96(%rsp,%rdx,4)
  jle .L6
  addl  $1, %eax
  movslq  %eax, %rdx
  movl  %ecx, 96(%rsp,%rdx,4)
  movslq  %esi, %rdx
  movl  %eax, 48(%rsp,%rdx,4)
  subl  $1, %esi
.L4:
  testl  %esi, %esi
  js .L12
  movslq  %esi, %rdx
  movl  (%rsp,%rdx,4), %ecx
  jmp .L5
.L12:
  ; (omitted)
  \end{lstlisting}
\end{frame}

\begin{frame}[fragile, shrink=30]{多重循环}
  \begin{lstlisting}[language=x86_64]
  movl  $1000000000, 96(%rsp) ; stk[0] = 1e9
  movl  $9, %esi              ; i = 9
  movl  $0, %eax              ; top = 0
  `\color{red}jmp .L4`
`\color{orange}.L6:`
  subl  $1, %eax              ; --top
`\color{cyan}.L5:`
  movslq  %eax, %rdx
  cmpl  %ecx, 96(%rsp,%rdx,4) ; x : stk[top]
  `\color{orange}jle .L6`                      ; jump if stk[top] <= x
  addl  $1, %eax              ; ++top
  movslq  %eax, %rdx
  movl  %ecx, 96(%rsp,%rdx,4) ; stk[top] = x
  movslq  %esi, %rdx
  movl  %eax, 48(%rsp,%rdx,4) ; lgis[i] = top (length of greedy LIS)
  subl  $1, %esi              ; --i
`\color{red}.L4:`
  testl  %esi, %esi
  js .L12                     ; NOT jump if i >= 0
  movslq  %esi, %rdx
  movl  (%rsp,%rdx,4), %ecx   ; x = ary[i]
  `\color{cyan}jmp .L5`
.L12:
  ; (omitted)
  \end{lstlisting}
\end{frame}

\section{switch 与跳转表}

\subsection*{C 中的 jump table}

\begin{frame}[fragile, shrink=40]{C 中的 jump table}
  \begin{lstlisting}{language=C}
switch (x) {
  case 0:
    x = 114;
  case 1:
    x += 514;
    break;
  case 3:
  case 4:
    x = 1919;
    break;
  default:
    x = 810;
}
  \end{lstlisting}

  $$
    \Updownarrow
  $$

  \begin{lstlisting}[language=C]
void bar(int& x) {
  void* jmpt[] = { &&loc_0, &&loc_1, &&loc_def, &&loc_3_4, &&loc_3_4 };
  if ((unsigned)x > 4) goto loc_def;
  goto *jmpt[x];
loc_0:
  x = 114;
loc_1:
  x += 514;
  goto done;
loc_3_4:
  x = 1919;
loc_def:
  x = 810;
done:;
}
  \end{lstlisting}
\end{frame}

\subsection*{汇编中的 jump table}

\begin{frame}[fragile, shrink=40]{汇编中的 jump table}
  \begin{lstlisting}[language=x86_64]
iterate_8per(int (*)(int), int, int):
  movl    %esi, %eax
  pushq   %r12
  movl    %esi, %r12d
  andl    $7, %eax
  pushq   %rbp
  pushq   %rbx
  movq    %rdi, %rbx
  movl    %edx, %edi
  jmp     *.L21(,%rax,8)
.L21:
  .quad   .L19
  .quad   .L27
  .quad   .L26
  .quad   .L25
  .quad   .L24
  .quad   .L23
  .quad   .L22
  .quad   .L20
.L20:
  call    *%rbx
  movl    %eax, %edi
.L22: ; two same lines
.L23: ; two same lines
.L24: ; two same lines
.L25: ; two same lines
.L26: ; two same lines
.L27: ; two same lines
  ; (omitted)
  \end{lstlisting}
\end{frame}

\begin{frame}[fragile, shrink=40]{汇编中的 jump table}
  \begin{lstlisting}[language=x86_64]
iterate_8per(int (*)(int), int, int):
  movl    %esi, %eax
  pushq   %r12
  movl    %esi, %r12d
  andl    $7, %eax
  pushq   %rbp
  pushq   %rbx
  movq    %rdi, %rbx
  movl    %edx, %edi
  `\color{red}jmp`     `\color{orange}*``\color{cyan}.L21`(,%rax,8)  ; switch here
`\color{cyan}.L21`:                     ; jump table for switch
  `\color{green}.quad   .L19`
  `\color{green}.quad   .L27`
  `\color{green}.quad   .L26`
  `\color{green}.quad   .L25`
  `\color{green}.quad   .L24`
  `\color{green}.quad   .L23`
  `\color{green}.quad   .L22`
  `\color{green}.quad   .L20`            ; targets for different cases
.L20:
  call    *%rbx
  movl    %eax, %edi
.L22: ; two same lines
.L23: ; two same lines
.L24: ; two same lines
.L25: ; two same lines
.L26: ; two same lines
.L27: ; two same lines
  ; (omitted)
  \end{lstlisting}
\end{frame}

\subsection*{引入 jump table 的优劣}

\begin{frame}{引入 jump table 的优劣}
  \begin{itemize}
    \item 优点: 
      \begin{itemize}
        \item 执行跳转指令的次数从 $\mathcal O(n)$ 降为 $\mathcal O(1)$, 适合 case 较多且分布较为密集的 switch.
      \end{itemize}
    \item 缺点:
      \begin{itemize}
        \item 需要额外的内存空间存储 jump table.
        \item case 分布过于稀疏时, jump table 会浪费大量空间.
        \item case 较少时, 直接使用条件测试 + 固定标签跳转反而更快.
      \end{itemize}
  \end{itemize}

  \pause
  \qquad \texttt{gcc} 会根据 case 数量和分布情况决定是否使用 jump table.\footnote{https://godbolt.org/z/8Y3zznaME}
\end{frame}

\subsection*{switch 语句的 if-else 翻译}

\begin{frame}{switch 语句的 if-else 翻译}
  \qquad 另一种 switch 的代码实现方式: if-else 链. 我们可以将所有 case 的内部语句依据 \texttt{break} 划分为若干个代码块, 显然程序最多进入一个代码块, 这可以用最外层的 if-else 链来控制.

  \,

  \pause
  \qquad 在一个代码块内, 由于没有 \texttt{break}, 程序会从某处开始顺序执行后续代码块的内容. 为此, 我们维护一个布尔变量, 记录 switch 变量是否已经触发过进入这个代码块的条件, 以决定代码块中的每条语句是否执行.
\end{frame}

\begin{frame}[fragile, shrink=5]{switch 语句的 if-else 翻译}
  \begin{lstlisting}[language=C]
  switch (x) {
    case 0: x += 5;
    case 1: x += 2;
    case 2: x *= 9; break;
    case 3: x *= 7;
    case 4: x *= 4;
    default: x -= 8;
  }
  \end{lstlisting}

  $$
    \Updownarrow
  $$

  \begin{lstlisting}[language=C++]
  int sx = x; // switch 内部语句可能修改 x
  bool entered = false;
  if (sx == 0 || sx == 1 || sx == 2) {
    if (entered |= sx == 0) x += 5;
    if (entered |= sx == 1) x += 2;
    x *= 9; // 已知进入此代码块, 所以一定会执行
  } else { // default 所在块: 无须再判断条件...
    if (entered |= sx == 3) x *= 7;
    if (entered |= sx == 4) x *= 4;
    x -= 8; // ...因为这句话一定需要执行
  }
  \end{lstlisting}
\end{frame}

\section{补充: 跳转的相关细节}

\subsection*{跳转指令的 PC 相对编码}

\begin{frame}[fragile,shrink=30]{跳转指令的 PC 相对编码}
  \begin{lstlisting}[language=x86_64]
0000000000000000 <popcount>:
   0: f3 0f 1e fa           endbr64 
   4: 55                    push   %rbp
   5: 48 89 e5              mov    %rsp,%rbp
   8: 89 7d ec              mov    %edi,-0x14(%rbp)
   b: c7 45 fc 00 00 00 00  movl   $0x0,-0x4(%rbp)
  12: eb `\color{red}13`                 `\color{orange}jmp`    27 <popcount+0x27>
  14: 8b 45 ec              mov    -0x14(%rbp),%eax
  17: 83 e0 01              and    $0x1,%eax
  1a: 89 c2                 mov    %eax,%edx
  1c: 8b 45 fc              mov    -0x4(%rbp),%eax
  1f: 01 d0                 add    %edx,%eax
  21: 89 45 fc              mov    %eax,-0x4(%rbp)
  24: d1 6d ec              shrl   -0x14(%rbp)
  27: 83 7d ec 00           cmpl   $0x0,-0x14(%rbp)
  2b: 75 `\color{red}e7`                 `\color{orange}jne`    14 <popcount+0x14>
  2d: 8b 45 fc              mov    -0x4(%rbp),%eax
  30: 5d                    pop    %rbp
  31: c3                    ret
  \end{lstlisting}

  \pause
  \qquad $\Longrightarrow$ 一般的跳转指令使用相对编码, 计算 jmp 的\textbf{下一条指令地址}与目标指令地址的差值, 将差值编码为一个 1/2/4 byte(s) 的\textbf{补码}.

  \pause
  \qquad \textbf{c.f.} 绝对地址跳转: \texttt{jmp  *\%rax}, \texttt{jmp  *(\%rax)}, ...
\end{frame}

\subsection*{使用条件传递消除跳转}

\begin{frame}{使用条件传递消除跳转}
  \qquad 跳转指令会破坏指令流水线, 影响 CPU 的指令级并行执行效率. 因此, 在某些情况下, 我们可以使用条件传递 (conditional move) 指令来消除跳转.

  \,

  \pause
  \qquad 这鼓励我们写出 "条件传递风格" 的代码\footnote{见 CS:APP p.379; 这里的目标和结论都偏经验.}.

  \,

  \pause
  \qquad $2\times$ faster 冒泡排序: https://godbolt.org/z/YTG3sEqTd (彩蛋: 在这个例子中, 你还将看到 a.笨笨编译器在函数末尾放两个 \texttt{ret} 用两个 label 跳转; b. \texttt{-O3} 默认的 slp-vectorize 昏招频出带来一倍负优化.)
\end{frame}

\subsection*{\&\&, || 与 goto}

\begin{frame}{\&\&, || 与 goto}
  \qquad \texttt{\&\&} 和 \texttt{||} 总是短路求值的, 这意味着在某个表达式求值时, 我们总能假定前面的子表达式已经求值完毕且均为真/假.

  \,

  \pause
  \qquad 这类似于若干个 if 语句的嵌套, 编译器一般也会将其翻译为这样的若干跳转指令.

  \,

  \pause
  \qquad 不过, 类似于编译器在已知求值简单时引入 \texttt{cmov} 这一优化, 编译器 (语义上) 也可能会将 \texttt{\&\&} 和 \texttt{||} 翻译不短路的 \texttt{\&} 和 \texttt{|}, 以减少跳转指令的使用.\footnote{https://godbolt.org/z/zMz4s3aTP}
\end{frame}

\begin{frame}[plain]
  % \vspace{0.4\textheight}
  \begin{center}
    \Huge\color{tyblue}\bfseries Thanks!
  \end{center}
\end{frame}

\end{document}
